% !PRINTED 534
\weberpage{534}
%<<CONT
wir ferner zwei entgegengesetzte Klassen durch $(A, \pm 2B, C)$, so sind
die entsprechenden Werte von $f(\omega)$ konjugiert imaginär, ihr Produkt
daher positiv, und es folgt also nach diesen Bestimmungen, daß das
Produkt
\[
\Pi f(\omega)
\]
einen positiven reellen Wert hat, der eine Potenz von 2 ist. Wir setzen,
indem wir mit $h$ die Klassenzahl bezeichnen, diese Potenz $= 2^{h\tau}$,
so daß
\weq{1}{\Pi\,\frac{f(\omega)}{2^\tau} = 1.}

Es wird unsere Aufgabe sein, diesen Exponenten $\tau$ zu bestimmen. Wir
schicken aber noch folgende Bemerkung voraus, die dieser Aufgabe ein
erhöhtes Interesse verleiht.

Infolge der Gleichung [§ 54, (8)]:
\weq{2}{f(\omega)^{24} - \gamma_2(\omega) f(\omega)^8 - 16 = 0}
ist, wenn $\omega$ die Wurzel einer zur Klasse $k$ gehörigen Form der
Diskriminante $-4m$ ist, $f(\omega)$ eine \gesperrt{ganze algebraische Zahl}, und da
$f(\omega)^8$ in (2) auch durch $-f_1(\omega)^8$, $-f_2(\omega)^8$ ersetzt
werden kann, so sind auch $f_1(\omega)$, $f_2(\omega)$ ganze algebraische
Zahlen. Mithin ist es auch
\[
\frac{\sqrt{2}}{f(\omega)} = f_1(\omega) f_2(\omega).
\]

Ist $p$ eine ungerade Primzahl, von der $-m$ quadratischer Rest ist, und
$p$ durch die Formen der Klasse $l$ (der Diskriminante $-4m$) darstellbar,
so ist bei passender Bestimmung von $c$ nach § 118:
\[
\frac{c+\omega}{p}
\]
die Wurzel einer zur komponierten Klasse $lk$ gehörigen Form, und es kann
(wenn nicht gerade $p=3$ ist), $c$ durch 48 teilbar angenommen werden.
Setzen wir also:
\[
u = f(\omega), \quad v = f\left(\frac{c+\omega}{p}\right),
\]
so ist sowohl $uv$ als $2:uv$ eine \gesperrt{ganze algebraische Zahl}.

Wenn wir daher nach § 74
\[
\begin{aligned}
B &= (uv)^s + \left(\frac{2}{p}\right)\frac{2^s}{(uv)^s},\\
A &= \left(\frac{u}{v}\right)^r + \left(\frac{v}{u}\right)^r
\end{aligned}
\]
%CONT>>
